\chapter{Methodology}
\label{ch:methodology}

\section{The Negotiation Cycle}

A single planning round proceeds as follows.

\begin{enumerate}[nosep]
  \item An order agent explodes its bill of materials into material and
        operation requirements.
  \item It announces its material requirement to the inventory agent and to
        every supplier agent, and its operation requirement to the relevant
        work-centre agents.
  \item The inventory agent offers the quantity it can reserve. Supplier
        agents offer the remainder at their own price and delivery date.
        Work-centre agents offer available slots.
  \item The order agent discards offers that cannot satisfy its constraints.
  \item It scores the remaining offers with deterministic, visible rules.
  \item It accepts the best valid combination; each acceptance becomes exactly
        one commitment.
  \item Resource agents record their commitments and reduce their available
        capacity or stock accordingly.
\end{enumerate}

\todo{sequence diagram of one full round as figures/negotiation-sequence.pdf}

\section{Bill of Materials Explosion}

\todo{state the recursive algorithm, with worked arithmetic for the
demonstration BOM. Include pseudocode in a listing}

\section{Inventory Netting and Reservation}

\todo{how gross requirement becomes net requirement, and how a reservation
prevents two orders from claiming the same stock}

\section{Offer Scoring}

Scoring must be deterministic and inspectable, because the ability to explain a
decision depends on it being reproducible.

\todo{give the scoring function explicitly: the factors, their weights, how
cost and lateness are normalised and traded off, and how ties are broken. State
the weights as a table so an examiner can recompute any score by hand}

\section{Commitment and Consistency}

\todo{how a commitment is made atomic, and what happens when two orders accept
overlapping offers from the same resource in the same round}

\section{Disruption Handling and Partial Replanning}

When a disruption is injected, the system identifies the commitments that
depended on the failed resource, releases only those, and starts a new
negotiation round for the affected unfinished work. Completed work and
commitments that did not depend on the failed resource are left untouched.

\todo{specify how the affected set is computed, and how the system proves it
did not disturb unaffected commitments --- this is the project's central claim
and needs to be measurable}

\section{Role of the Language Model}

The language model is confined to two tasks: parsing a free-text draft order
into a structured order, and producing read-only explanations of a recorded
negotiation. It never performs arithmetic, never judges feasibility, never
scores an offer, and never writes state. Parsed output is validated against a
schema and rejected if it does not conform, so a bad generation becomes a
visible failure rather than a silently wrong plan.

\todo{show the schema and the validation path, and describe the degraded
behaviour when no model is reachable}
